% SPDX-License-Identifier: Apache-2.0
% covers: FrameIn
\datasheet{FrameIn}{The Ethernet receiver's bytes packed into words for
a store engine, with the frame's length and the slot it lands in.}

\dsfacts{\code{lib/parts/src/ethdma.rs}}{\code{txhdl\_parts::ethdma::FrameIn}}%
{\code{//docs:examples}}

\subsection*{Function}

The mirror of \code{FrameOut}. It takes a frame's bytes on
\code{rx} and offers them to \code{LineStore} as 32-bit words on
\code{words}, lowest lane first. It does not touch the bus and does
not know where the frame lands in memory.

It needs the frame's length on \code{len} before the first byte,
because the store engine is told how many bytes to write when it
starts. \code{EthRx} offers exactly that, since it takes and checks a
whole frame before offering any of it, and so does \code{FrameLen}.

It also says three things to the register block: \code{count}, the
length of the frame taken; \code{store}, high while a frame is being
taken, which starts the store engine; and \code{which}, the slot the
frame lands in, which alternates.

\subsection*{Parameters}

None. The length is 16 bits wide, the words 32, and there are two
slots.

\dstables{FrameIn}

\subsection*{Behaviour}

\begin{itemize}
\item A frame starts when a byte is offered, \code{len} is not zero,
and \code{hold} is low. The length is latched there, and the slot
moves to the other one.
\item Bytes are packed from the top of the word down, so after four
the first is in the low lane.
\item The last word of a frame is usually partial, and it is sent with
its real bytes shifted into the low lanes, where the store engine's
strobe expects them. The engine strobes the rest away, so nothing past
the frame is written.
\item \textbf{\code{hold} keeps it from starting the next frame while
the store engine is busy with the one before.} This unit finishes
handing over words well before the engine's write is answered, and the
register block reads \code{count} and \code{which} when the engine
goes idle. A new frame started early would change both under the
frame still being stored, and the driver would be told the new frame's
length and slot for the old frame's bytes. Only a run with a real
store engine behind this unit can see that; it was found by reading
the board join, not by a test.
\end{itemize}

\subsection*{Verification}

\begin{itemize}
\item \code{ex\_ethdma} is \code{FrameOut}, the real \code{EthTx} and
\code{EthRx}, and this unit in a line. Two frames of 101 and 67 bytes
go through it, and every word that lands is checked, the partial last
words included. The run also checks that four words past the second
frame are untouched, that the receiver's length for the last frame
arrived, and that the second frame took the other slot.
\item The netlist is checked against that run under nvc and
Verilator: \code{//docs:ethdma\_sim\_frame\_in\_frame\_in\_tb\_test}
and \code{//docs:ethdma\_vsim\_frame\_in\_test}.
\item \code{ex\_framelen} runs it behind \code{FrameLen} instead of
\code{EthRx}, with the same two frames.
\item On the board,
\code{a\_frame\_received\_lands\_in\_memory\_and\_the\_core\_reads\_it\_back}
in \code{cpu/vreteno/tests/board.rs} injects a frame on the wire and
has the core read it back from its slot.
\end{itemize}

\subsection*{Used by}

\code{Board}, as \code{fin}, behind \code{FrameLen} and in front of the
store engine \code{store}. \code{hold} is the store engine's
\code{running}, and \code{count} and \code{which} go to
\code{EthSlots} as \code{rx\_len} and \code{rx\_which}. The examples
\code{ex\_ethdma} and \code{ex\_framelen}.

\subsection*{Limits}

Its ports are named apart from \code{FrameOut}'s on purpose. A
testbench reads the trace by port name, so two lowered units in one
run that both name a port \code{out} bind to the same recorded signal
(issue 462).

A frame whose length is zero is never started, and its bytes stay in
the channel.

It runs on one clock, \code{DefaultClock}.
